% Part: cheat-sheets
% Chapter: quantified-modal-logics
% Section: syntax

\documentclass[../../../../include/open-logic-section]{subfiles}

\begin{document}

\olfileid{chs}{qml}{syn}

\olsection{Syntax}

Quantified modal logic combines the languages of
quantification (first-order logic with identity, free logic) and of
modal logic. We add an existence predicate, as in free logic. 

\begin{enumerate}
\item Logical symbols
\begin{enumerate}
\item Logical connectives:
  \startycommalist
  \iftag{prvFalse}{\ycomma $\lfalse$ (the propositional constant for !!{falsity})}{}%
  \iftag{prvTrue}{\ycomma $\ltrue$ (the propositional constant for !!{truth})}{}%
  \iftag{prvNot}{\ycomma $\lnot$ (negation)}{}%
  \iftag{prvAnd}{\ycomma $\land$ (conjunction)}{}%
  \iftag{prvOr}{\ycomma $\lor$ (disjunction)}{}%
  \iftag{prvIf}{\ycomma $\lif$ (!!{conditional})}{}%
  \iftag{prvIff}{\ycomma $\liff$ (!!{biconditional})}{}%
  \iftag{prvAll}{\ycomma $\lforall$ (universal quantifier)}{}%
  \iftag{prvEx}{\ycomma $\lexists$ (existential quantifier)}{}%
  \iftag{prvBox}{\ycomma $\Box$ (necessity operator)}{}%
  \iftag{prvDiamond}{\ycomma $\Diamond$ (possibility operator)}{}.
\item The one-place existence predicate~$\lfrexists$.
\item The two-place !!{identity}~$\eq$.
\item A !!{denumerable}s set of !!{variable}s: $\Obj v_0$, $\Obj v_1$, $\Obj
  v_2$, \dots
\end{enumerate}
\item Non-logical symbols, making up the \emph{standard
  language} of quantified modal logic
\begin{enumerate}
\item A !!{denumerable}s set of $n$-place !!{predicate}s for each $n>0$: $\Obj
  A^n_0$, $\Obj A^n_1$, $\Obj A^n_2$, \dots
\item A !!{denumerable}s set of !!{constant}s: $\Obj c_0$, $\Obj c_1$, $\Obj
  c_2$, \dots.
\item A !!{denumerable}s set of $n$-place !!{function}s for each $n>0$:
  $\Obj f^n_0$, $\Obj f^n_1$, $\Obj f^n_2$, \dots
\end{enumerate}
\item Punctuation marks: (, ), and the comma.
\end{enumerate}

A language~$\Lang L$ for quantified modal logic is 
build from a selection of all or some
!!{predicate}s, !!{constant}s, and !!{function}s as follows.

\begin{defn}[Terms]
\ollabel{defn:terms}
The set of \emph{terms}~$\Trm[L]$ of~$\Lang L$ is
defined inductively by:
\begin{enumerate}
\item Every !!{variable} is a term.
\item Every !!{constant} of~$\Lang L$ is a term.
\item If $f$ is an $n$-place !!{function} and $t_1$, \dots, $t_n$
  are terms, then $\Atom{f}{t_1, \ldots, t_n}$ is a term.
\tagitem{limitClause}{
  Nothing else is a term.}{}
\end{enumerate}
A term containing no !!{variable}s is a \emph{closed term}.
\end{defn}

\begin{defn}[Formulas]
\ollabel{defn:formulas}
The set of \emph{!!{formula}s}~$\Frm[L]$ of the language~$\Lang L$
is defined inductively as follows:
\begin{enumerate}
\tagitem{prvFalse}{$\lfalse$ is an atomic !!{formula}.}{}

\tagitem{prvTrue}{$\ltrue$ is an atomic !!{formula}.}{}

\item If $R$ is an $n$-place !!{predicate} of~$\Lang L$ and $t_1$, \dots,
  $t_n$ are terms of~$\Lang L$, then $\Atom{R}{t_1,\ldots, t_n}$ is an
  atomic !!{formula}.

\item If $t$ is a term of~$\Lang L$, then $\Atom{\lfrexists}{t}$
  is an atomic !!{formula}.

\item If $t_1$ and $t_2$ are terms of~$\Lang L$, then $\Atom{\eq}{t_1, t_2}$
  is an atomic !!{formula}.

\tagitem{prvNot}{If $!A$ is !!a{formula}, then $\lnot !A$ is
  !!a{formula}.}{}

\tagitem{prvAnd}{If $!A$ and $!B$ are !!{formula}s, then $(!A \land
  !B)$ is !!a{formula}.}{}

\tagitem{prvOr}{If $!A$ and $!B$ are !!{formula}s, then $(!A \lor !B)$
  is !!a{formula}.}{}

\tagitem{prvIf}{If $!A$ and $!B$ are !!{formula}s, then $(!A \lif !B)$
  is !!a{formula}.}{}

\tagitem{prvIff}{If $!A$ and $!B$ are !!{formula}s, then $(!A \liff !B)$
  is !!a{formula}.}{}

\tagitem{prvAll}{If $!A$ is !!a{formula} and $x$ is !!a{variable},
  then $\lforall[x][!A]$ is !!a{formula}.}{}

\tagitem{prvEx}{If $!A$ is !!a{formula} and $x$ is !!a{variable},
  then $\lexists[x][!A]$ is !!a{formula}.}{}

\tagitem{prvBox}{If $!A$ is !!a{formula} then $\Box !A$ is
!!a{formula}.}{}

\tagitem{prvDiamond}{If $!A$ is !!a{formula} then $\Diamond !A$ is !!a{formula}.}{}

\tagitem{limitClause}{Nothing else is !!a{formula}.}{}
\end{enumerate}
\end{defn}

Free and bound !!{variable}s, !!{sentence}s, variable substitution are defined as in
first-order logic.

\Article{sentence} \emph{!!{sentence}} is !!a{formula} without free
!!{variable}s.

\end{document}